\begin{flushright}
        \scriptsize{Не бывает безвыходных ситуаций. \\ Бывают ситуации, выход из которых тебя не устраивает. \\
        Нара Шикамару}
\end{flushright}


\section{Исследование 1}

\problem{1} Расставьте все числа от 1 до 9 в квадрате $3 \times 3$ таким образом, что  соседние числа имеют разные четности.

\problem{2}
В шеренгу выстроились 100 рыцарей и лжецов. Первый сказал: «Количество рыцарей среди нас — делитель числа 1», второй сказал: «Количество рыцарей среди нас — делитель числа 2», и так далее. Сколько в шеренге рыцарей?
% 10

\problem{3} 
В ряд выписаны числа от 1 до 10. Можно ли между ними расставить знаки $ + $ и $ - $ так, чтобы получился 0?

% Нельзя, так как среди чисел от 1 до 10 нечётное количество нечётных.

\problem{4}
На прогулку пошли шестиклассники и пятиклассники. Все они были либо босиком, либо в тапочках. Шестиклассников было 24, а босых учеников 16. Обутых пятиклассников было столько же, сколько босых шестиклассников. Сколько учеников ходили на прогулку?
% 40


\problem{5}
В особой больнице есть главврач и много пациентов. В течение недели каждый пациент один раз в день кусал кого-нибудь (возможно и себя). В конце недели оказалось, что у всех пациентов по 2 укуса, а у главврача — 100 укусов. Сколько в больнице пациентов?
% х = 20

\newpage

\section{Исследование 2}

\problem{1} В коробке 14 белых и 14 чёрных шариков. Какое минимальное количество шариков нужно достать из коробки, чтобы среди них наверняка оказалось 2 черных шарика?
% 16


\problem{2} Представте 49 как сумму 20 слагаемых с одинаковой суммой цифр.
% 11 + 2 + 2 и тд

\problem{3}Сколько раз за сутки на часах минутная стрелка обгонит часовую?
% 21


\problem{4} В классе 33 ученика. 27 из них изучают английский язык, и 10 — французский. Сколько учеников могут изучать и английский, и французский? 
% 4

\problem{5} Можно ли разложить 44 шарика на 9 кучек так,
чтобы количество шариков во всех кучках было различным (и ненулевым)?
% нельзя

\newpage

\section{Исследование 3}

\problem{1} Произведение двух чисел умножили на их разность. Могло ли получиться 30?
% могло (5-3) * 15

\problem{2} В кружочках написана цифры и знаки арифметических действий, как показано на рисунке справа. Разрешается переходить из любого кружочка в соседний, записы- вая при этом в строку то, что написано в кружочках. Как, начав с одного из них, обойти всё по одному разу так, чтобы записанное равенство оказалось верным?
\begin{center}
\includegraphics[width=0.4\textwidth]{Летние Сборы/Сборы Элиста 2024/ВВ/лабиринт_3_1.png}
\end{center}

\problem{3} Взяли 17 целых чисел и сложили. Результат суммы поделили на 17 и получили 63.5.  Могло ли такое получиться? 
% не может, из-за того что сумма целая


\problem{4} Обезьяна становится счастливой, когда съедает три разных фрукта. Какое наибольшее количество обезьян можно осчастливить, имея 20 груш, 30 бананов, 40 персиков и 50 мандаринов?
% 45

\problem{5}
На этот раз амеба Лунтик исправила две цифры в примере на умножение. Получилось $4 * 5 * 4 * 5 * 4 = 2247$. Помогите профессору Составителю восстановить исходный пример.
% 4 * 5 * 4 * 7 * 4 = 2240

